\documentclass{article}
\usepackage[T1]{fontenc}
\usepackage[utf8]{inputenc}
\usepackage{amsmath,amssymb}

\begin{document}

\section{Simulation study on $\mathbb{S}^1$}

We evaluate frequentist risk performance of the proposed denoising estimators on the unit circle ($\mathbb{S}^1$). Let $\Theta \in \mathbb{S}^1$ denote a latent signal and $X \in \mathbb{S}^1$ its noisy observation generated from the isotropic Riemannian normal model
\[
X \mid \Theta \sim \mathcal{N}_{\mathbb{S}^1}\!\left(\Theta,\sigma^2\right),
\]
implemented via \texttt{random\_riemannian\_normal} available in \texttt{Geomstats}. Performance of a denoiser $\delta$ is measured using squared geodesic loss,
\[
\mathcal{R}(\delta) := \mathbb{E}\!\left[d^2\!\left(\delta(X),\Theta\right)\right],
\]
where $d(\cdot,\cdot)$ is the geodesic distance on $\mathbb{S}^1$.

To probe robustness across increasing multimodality, we consider four sampling schemes $G$ for $\Theta$,
namely: (i) one--modal, a single Riemannian Gaussian mode with variance $\tau^2=0.05$; (ii) two--modal, a two-component Riemannian Gaussian mixture with separated modes and component variance $\tau^2=0.04$; (iii) three--modal, a three-component separated mixture with component variance $\tau^2=0.03$; and (iv) four--modal, a four-component separated mixture with component variance $\tau^2=0.02$.
For visualization we display a circular kernel density estimate of each $G$ on a polar plot.

We compare the Na\"{\i}ve baseline $\delta_N(X)=X$, corresponding to no denoising, with the oracle Tweedie denoiser $\delta_T$ and the oracle Bayes estimator $\delta_B$, which in the (rather exceptional) case of the circle can be computed in closed form. These oracle procedures are computed by Monte Carlo integration using independent sampling from the true $G$ and therefore serve as benchmarks for the best achievable behavior under the model (bypassing density/score estimation). Concretely, $\delta_T$ is based on approximating the conditional expectation $\mathbb{E}[\log_X \Theta \mid X]$, while $\delta_B$ approximates the Bayes rule by rejection sampling.

We first examine how risk varies with the noise variance $\sigma^2$. For each target $G$, we evaluate the Monte Carlo risk
\[
\widehat{\mathcal{R}}(\delta) := \frac{1}{n}\sum_{i=1}^n d^2\!\left(\delta(X_i),\Theta_i\right)
\]
with $n=500$. We consider increasing values of $\sigma^2$ on a grid in $[0.005,0.25]$ and repeat the experiment across $N_{\mathrm{MC}}=3$ Monte Carlo replications. The output is plotted with mean risk and 68\% confidence intervals across replications.

This experiment highlights the dependence of denoising gains on the underlying signal distribution $G$. As expected, for unimodal targets the oracle procedures substantially improve over $\delta_N$ as $\sigma^2$ increases. As multimodality increases, the gap between $\delta_T$ and $\delta_B$ becomes more visible at moderate-to-large noise levels, reflecting the increasing difficulty of denoising under highly non-uniform, multi-peak structure.

The second experiment studies the finite-sample excess risk induced by estimating the denoiser from data relative to its oracle counterpart. Concretely, for each $G$, each $\sigma^2$, and each training sample size $N \in \{50,100,250,500\}$, we draw an independent test set, compute the denoiser $\hat\delta_T(\cdot;\rho)$ using fixed complexity ($M=7$) and a grid of regularization parameters $\rho$, and select $\rho^*$ by minimizing the empirical test risk. We then compute the oracle Tweedie benchmark on the same test set and display the difference
\[
\Delta(N;\sigma^2,G)
:= \widehat{\mathcal{R}}_{\mathrm{test}}\!\left(\hat\delta_T(\cdot;\rho^*)\right)
- \widehat{\mathcal{R}}_{\mathrm{test}}(\delta_T)
\]
across $N_{\mathrm{MC}}=10$ replications. We plot $\Delta$ versus $N$ on a log ($x$)-scale, with separate curves for each $\sigma^2$.

The excess risk decays with $N$, consistent with improved stability of the estimated score/denoiser and more reliable tuning as the training sample size increases. The decay is slower and the gap larger for higher $\sigma^2$ and for more multimodal targets, where accurate estimation is more challenging.

In summary, the $\mathbb{S}^1$ simulations demonstrate: (i) substantial oracle gains over the Na\"{\i}ve estimator when $G$ is concentrated; (ii) increasing multimodality makes denoising harder, amplifying the gap between $\delta_T$ and $\delta_B$ at moderate-to-large noise; and (iii) the excess risk of empirically tuned procedures over the oracle benchmark decreases with training sample size but is most pronounced for high noise and highly multimodal regimes.
% ...existing code...


\section*{Simulation study on $\mathbb{S}^2$}
We evaluate frequentist risk performance of the proposed denoising estimators on the two--dimensional sphere ($\mathbb{S}^2$). Let $\Theta \in \mathbb{S}^2$ denote a latent signal and $X \in \mathbb{S}^2$ its noisy observation generated from the isotropic Riemannian normal model
\[
X \mid \Theta \sim \mathcal{N}_{\mathbb{S}^2}\!\left(\Theta,\sigma^2\right),
\]
implemented via \texttt{random\_riemannian\_normal} available on \texttt{Geomstats} python module. Performance of a denoiser $\delta$ is measured using squared geodesic loss, i.e.
\[
\mathcal{R}(\delta) := \mathbb{E}\!\left[d^2\!\left(\delta(X),\Theta\right)\right],
\]
where $d(\cdot,\cdot)$ is the geodesic distance on $\mathbb{S}^2$.

To probe robustness across qualitatively different signal structures, we consider four sampling schemes $G$ for $\Theta$, each implemented through \texttt{get\_G\_class('S2', sampler, name, params)}:
\begin{itemize}
\item Uniform: $\Theta\sim \mathrm{Unif}(\mathbb{S}^2)$.
\item One--modal: a single concentrated mode, i.e.\ a Gaussian Riemannian sampler with variance $\tau^2=0.05$.
\item Four--modal: a Gaussian Riemannian mixture sampler with four separated modes ($\tau^2=0.01$ for each).
\item Equatorial: mass concentrated around the equator geodesic.
\end{itemize}

These four targets span regimes from fully diffuse (uniform) to strongly non-uniform and highly structured (multimodal/equatorial). For visualization we display a kernel density estimate of each $G$ on a Mollweide projection.

We a Na\"{\i}ve baseline $\delta_N(X)=X$, corresponding to no denoising,
with the oracle denoiser $\delta_T$.

The latter is ``oracle'' in the sense that it have access to independent sampling from the true $G$ and therefore serve as benchmarks for the best achievable behavior under the model, and does not go thrugh estimation of the density/score, but rather by the (approximat) conditional expecation $\mathbb{E}[\log_X\Theta \mid X]$ computed by Monte Carlo integration.


We first examine how risk varies with the noise variance $\sigma^2$. For each target $G$, we evaluate the monte carlo risk
\[
\widehat{\mathcal{R}}(\delta) := \frac{1}{n}\sum_{i=1}^n d^2\!\left(\delta(X_i),\Theta_i\right)
\]
for some large $n$, say $n=500$.
We consider increasing values of$\sigma^2 \in \{0.005,0.01,0.05,0.10, 0.15, 0.20, 0.25\}$  and repeat the experiment across $N_{\mathrm{MC}}=20$ Monte Carlo replications. The output is plotted with mean risk and 68\% confidence intervals across replications.

This experiment quantifies the expected ordering between the no-denoising baseline $\delta_N$ and the oracle procedure as noise increases, and highlights the dependence of the potential denoising gains on the underlying signal distribution $G$. As expected, under the uniform target (no structure to exploit) the na\"{\i}ve and oracle denoisers are nearly indistinguishable, whereas for structured targets (modal or equatorial) the oracle denoisers leverage concentration and exhibit substantially smaller risk, especially at moderate-to-large $\sigma^2$.

The second experiment studies the finite-sample excess risk induced by estimating the denoiser from data and its oracle version. Concretely, for each $G$, each $\sigma^2$, and each training sample size $N \in \{50,100,250,500,1000\}$, we
draw an independent sample $\{\Theta_i\}_{i=1}^{N}$ and noisy $X_i$; wecompute the denoiser $\hat\delta_T(\cdot;\rho)$ using  with fixed complexity ($M=7$) and cross-validated regularisation $\rho$;  we compute the oracle and empirical  Tweedie risk on the sets, and display the Monte Carlo distribution of the difference
\[
\Delta(N;\sigma^2,G)
:= \widehat{\mathcal{R}}_{\mathrm{test}}\!\left(\hat\delta_T(\cdot;\rho^*)\right)
- \widehat{\mathcal{R}}_{\mathrm{test}}(\delta_T),
\]
 between the two risks across $20$ replications. We plot $\Delta$ versus $N$ on a log ($x$)-scale, with separate curves for each $\sigma^2$.

 The curves decay with $N$, consistently with the theory, since larger  sample size improves density/score estimation and stabilizes tuning


 In summary, the $\mathbb{S}^2$ simulations demonstrate: (i) oracle denoisers can yield substantial risk reductions relative to $\delta_N$ when $G$ is structured; (ii) the excess risk of empirically tuned procedures over their oracle counterpart decreases with training sample size and is most visible in high-noise, structured regimes.
%   and evaluate the
%  test risk
% \[
% \widehat{\mathcal{R}}_{\mathrm{test}}\!\left(\hat\delta_T(\cdot;\rho)\right)
% =\frac{1}{m}\sum_{j=1}^{m} d^2\!\left(\hat\delta_T(X_j^{\mathrm{test}};\rho),\Theta_j^{\mathrm{test}}\right).
% \]
% \item Select $\rho^*$ minimizing the Monte Carlo average of $\widehat{\mathcal{R}}_{\mathrm{test}}$ over replications.
% For each configuration, we record the Monte Carlo distribution of the difference
% \[
% \Delta(N;\sigma^2,G)
% := \widehat{\mathcal{R}}_{\mathrm{test}}\!\left(\hat\delta_T(\cdot;\rho^*)\right)
% - \widehat{\mathcal{R}}_{\mathrm{test}}(\delta_T),
% \]
% and summarize it by the mean, median, and standard deviation across $N_{\mathrm{MC}}=10$ replications (object \texttt{df\_N}).


\end{document}